\phantomsection\addcontentsline{toc}{section}{High-Speed Transportation}
\ECchapter{30}{High-Speed Transportation}{Optimising speed, energy, capacity and safety}{ECBlue}

\ECObjectives{
\begin{itemize}
\item Explain how aerodynamic drag, propulsion power and braking distance scale with speed.
\item Compare conventional rail, magnetic levitation and low-pressure concepts at system level.
\item Estimate stopping distance and define energy, infrastructure and safety requirements for a corridor.
\end{itemize}}

\begin{multicols}{2}
\begin{engineeringnews}
High-speed rail, magnetic-levitation systems and low-pressure-tube concepts seek to reduce journey times between
cities. At high speed, aerodynamic resistance, guideway accuracy, signalling, traction power and emergency
procedures become dominant constraints. The vehicle and infrastructure must therefore be designed as one system.
\end{engineeringnews}

\begin{ecarticle}
High-speed transportation is not simply a conventional vehicle with a larger motor. Track geometry, propulsion,
power supply, signalling, stations, maintenance and evacuation must be designed together. At $300\,\mathrm{km\,h^{-1}}$,
a train travels at $83.3\,\mathrm{m\,s^{-1}}$, so a five-second delay adds more than four hundred metres to stopping
distance before effective braking even begins.

Aerodynamic drag can be approximated by
\[
F_D=\frac{1}{2}\rho C_D A v^2,
\]
where $\rho$ is air density, $C_D$ the drag coefficient, $A$ frontal area and $v$ speed. The associated power is
\[
P_D=F_Dv=\frac{1}{2}\rho C_D A v^3.
\]
Thus, doubling speed multiplies aerodynamic power by approximately eight if geometry and air density remain
unchanged. Streamlined noses, sealed inter-car gaps and protected underbodies reduce drag and pressure fluctuations.
Tunnel entry is especially demanding because the train compresses air in a confined space.

Steel-wheel high-speed rail offers low rolling resistance, mature switching technology and partial compatibility
with existing networks. Magnetic levitation removes wheel--rail contact at speed and can reduce mechanical wear,
but requires dedicated guideways, accurate air-gap control and specialised power electronics. Low-pressure tubes
could reduce aerodynamic loss further, although maintaining a sealed structure, controlling thermal expansion and
providing safe evacuation are major unresolved system constraints.

For constant deceleration $a$, ideal braking distance is
\[
d_b=\frac{v^2}{2a}.
\]
Real stopping distance also contains sensing, communication, decision and brake-build-up delays. High-speed lines
therefore use continuous cab signalling and automatic train protection. Safe separation must include uncertainty in
adhesion, gradient, train mass, actuator response and localisation.

Regenerative braking can return electrical energy when the supply network, another train or a storage system can
absorb it. Recovery is limited by conversion efficiency, timetable coincidence, grid receptivity and auxiliary
loads. An energy balance must also include rolling resistance, gradients, acceleration, HVAC and station operation.

Maximum speed alone is a poor indicator of usefulness. Door-to-door time includes access, waiting, dwell and
transfer times. Capacity, reliability, noise, construction carbon and passenger demand can dominate the lifecycle
assessment. A high-speed corridor is justified only when system benefits exceed the infrastructure, energy and
maintenance costs over decades of operation.
\end{ecarticle}

\begin{tcolorbox}[title=\sffamily\bfseries Engineering insight,colback=yellow!10,colframe=ECOrange]
A higher maximum speed may save little time when stations are close together. Acceleration, braking and dwell time
must be included before selecting the design speed.
\end{tcolorbox}

\begin{engineeringfocus}
\textbf{System-level energy and safety chain}
\begin{center}
\resizebox{0.96\linewidth}{!}{%
\begin{tikzpicture}[font=\scriptsize,>=Latex,node distance=3.5mm]
\node[draw=ECBlue,fill=ECLightBlue,rounded corners] (grid) {power supply};
\node[draw=ECPurple,fill=ECPurple!10,rounded corners,right=of grid] (drive) {converter and motor};
\node[draw=ECGreen,fill=ECGreen!10,rounded corners,right=of drive] (motion) {vehicle motion};
\node[draw=ECOrange,fill=ECOrange!10,rounded corners,below=of drive] (loss) {drag and auxiliaries};
\node[draw=ECRed,fill=ECRed!8,rounded corners,below=of motion] (safe) {signalling and braking};
\draw[->,thick] (grid)--(drive); \draw[->,thick] (drive)--(motion);
\draw[->,thick] (motion)--(loss); \draw[->,thick,dashed] (safe)--(motion);
\draw[->,thick,dashed] (motion) to[bend left=35] node[above]{regeneration} (grid);
\end{tikzpicture}}
\end{center}
\end{engineeringfocus}

\begin{keyfigures}
\begin{tabularx}{\linewidth}{@{}Xr@{}}
Drag force & $\propto v^2$\\
Aerodynamic power & $\propto v^3$\\
Ideal braking distance & $v^2/(2a)$\\
Safety architecture & continuous protection\\
Useful metric & passenger-kilometre
\end{tabularx}
\end{keyfigures}

\begin{tcolorbox}[title=\sffamily\bfseries Aerodynamic power versus speed,colback=white,colframe=ECBlue]
\begin{center}
\begin{tikzpicture}
\begin{axis}[width=0.94\linewidth,height=45mm,xlabel={Speed (km/h)},ylabel={Relative aerodynamic power},
grid=major,ticklabel style={font=\scriptsize},label style={font=\scriptsize}]
\addplot[thick,mark=*] table[x=speed,y=power,col sep=comma]{data/EC30/power.csv};
\end{axis}
\end{tikzpicture}
\end{center}
\footnotesize The cubic trend explains why modest increases in cruise speed can require large increases in traction energy.
\end{tcolorbox}

\begin{applicationexercise}{Stopping-distance and power budget}
A train has $C_DA=8.0\,\mathrm{m^2}$ and runs in air of density $1.2\,\mathrm{kg\,m^{-3}}$ at
$300\,\mathrm{km\,h^{-1}}$.
\begin{enumerate}
\item Convert the speed to $\mathrm{m\,s^{-1}}$ and estimate aerodynamic drag and power.
\item With emergency deceleration $0.9\,\mathrm{m\,s^{-2}}$, calculate the ideal braking distance.
\item Add a five-second response delay and determine the total stopping distance.
\item Apply a 15\% safety margin and propose a minimum protected distance.
\item Estimate the aerodynamic power at $330\,\mathrm{km\,h^{-1}}$ using the cubic scaling law and discuss the time--energy trade-off.
\end{enumerate}
\end{applicationexercise}

\begin{vocabulary}
\begin{tabularx}{\linewidth}{@{}lX@{}}
aerodynamic drag & traînée aérodynamique\\
rolling resistance & résistance au roulement\\
guideway & voie de guidage\\
magnetic levitation & sustentation magnétique\\
pressure wave & onde de pression\\
braking distance & distance de freinage\\
regenerative braking & freinage régénératif\\
train protection & protection automatique des trains\\
passenger-kilometre & passager-kilomètre\\
dedicated infrastructure & infrastructure dédiée
\end{tabularx}
\end{vocabulary}

\begin{questions}
\item Why does aerodynamic power rise more quickly than drag force?
\item Compare the main advantages and limitations of magnetic levitation.
\item Why are continuous signalling and automatic protection necessary?
\item Under what conditions can regenerative braking recover useful energy?
\item Why can maximum speed be a poor measure of door-to-door performance?
\end{questions}

\begin{engineeringchallenge}
Compare a $300\,\mathrm{km\,h^{-1}}$ railway and a magnetic-levitation line for a 600 km corridor. Build a requirement
table covering journey time, capacity, energy, protected stopping distance, guideway tolerance, evacuation and
availability. Estimate non-stop travel time and aerodynamic power, identify three degraded modes, and propose a
lifecycle decision matrix including infrastructure carbon and passenger demand.
\end{engineeringchallenge}

\begin{speaking}
Should public investment favour maximum speed or a denser network with lower speeds and more frequent services?
Discuss accessibility, capacity, energy, regional equality and lifecycle cost.
\end{speaking}
\end{multicols}

\subsection*{Sources}
\ECsource{International Union of Railways -- High-speed rail}{https://uic.org/passenger/highspeed/}
\ECsource{International Energy Agency -- Rail}{https://www.iea.org/energy-system/transport/rail}